\section{Instrução ao VASP}

\subsection{Funcionalidades dos arquivos de entrada}
O VASP utiliza arquivos de entrada específicos para definir os parâmetros da simulação. Os
principais arquivos de entrada são:

\begin{itemize}
    \item \texttt{INCAR}: Contém os parâmetros de controle da simulação;
    \item \texttt{POSCAR}: Define a estrutura cristalina e as posições dos átomos;
    \item \texttt{KPOINTS}: Especifica a malha de pontos; e
    \item \texttt{POTCAR}: Contém os pseudopotenciais dos elementos presentes na simulação.
\end{itemize}

\subsubsection{Arquivo \texttt{INCAR}}
O arquivo \texttt{INCAR} é um arquivo de \textit{input} central no VASP, que determina \emph{o que
fazer e como fazer} durante a simulação. As \textit{tags} presentes no arquivo \texttt{INCAR} são o
conjunto de parâmetros de controle que definem o tipo de cálculo a ser realizado e os algoritmos a
serem utilizados pelo VASP. Por ser uma das principais causas de erros durante a execução do VASP, é
importante compreender o significado de cada \textit{tag} e como elas afetam o comportamento do código.

O formato de cada instrução consiste no nome de uma \textit{tag}, no sinal de igual (=) e nos valores
atribuídos à \textit{tag}, por exemplo, (\texttt{tag = valores}). Segue abaixo o arquivo
\texttt{INCAR} utilizado para a simulação do EG-Graphene, com comentários explicativos sobre cada \textit{tag}:

\begin{minted}[linenos, frame=single, bgcolor=lightgray!20]{fortran}
Electronic relaxation:

ENCUT   = 800.000    !energia de cutoff
ALGO    = Normal     !algoritmo para minimizacao eletronica
NELM    = 500        !valor default para numero de passos SCF
NELMIN  = 6          !numero minimo de passos SCF
NELMDL  = -12        !numero de passos NSCF no inicio
EDIFF   = 1.0E-7     !criterio de convergencia SCF
AMIX    = 0.1        !mix densidade pmix=(1-AMIX)*pin+AMIX*pout
BMIX    = 0.0001     !outra variavel para misturar densidade

Calculation mode :

PREC    = Accurate   !influencia a FFT
IVDW    = 0          !interacao vdW Grimme D3
ISPIN   = 1          !spin nao polarizado
ADDGRID = .TRUE.     !grid add para diminuir ruido no cal das forcas
LASPH   = .TRUE.

Ionic relaxation:

NSW     = 200        !numero maximo de passos na relaxacao
EDIFFG  = -0.010     !para a relaxacao quando as forcas forem menores que
                     !|EDIFFG|
IBRION  = 2          !algoritmo para relaxacao da rede, POTIM não interfere
ISIF    = 3          !relaxacao da celula unitaria e ions, muda forma e
                     !volume da celula (2 pode vir a ser melhor)
POTIM   = 0.5
LATTICE_CONSTRAINTS = .TRUE. .TRUE. .FALSE

Integration over the Brillouin zone (BZ):

ISMEAR  = 0           !Gaussian-smearing
SIGMA   = 0.01        !largura do smearing em eV

DOS calculation:

LORBIT  = 10          !escreve DOSCAR e PROCAR
NEDOS   = 6001        !numero de pontos no grid da DOS

OUTCAR size:

NWRITE  = 1           !informacao escrita no OUTCAR
LWAVE   = .FALSE.     !nao escreve o WAVECAR
LCHARG  = .FALSE.     !nao escreve o CHGCAR e CHG

Key for parallel mode calculation:

NCORE   = 4           !numero de cores por orbital
LPLANE  = .TRUE.
\end{minted}

Analisando o meu arquivo de entrada \texttt{INCAR}, o primeiro parâmetro é o \texttt{ENCUT}, que
define a energia de corte para a base de ondas planas em $\unit{\electronvolt}$, sendo o valor de
$\SI{800}{\eV}$ utilizado para garantir a convergência da energia total do sistema.

\subsubsection{Arquivo \texttt{POSCAR}}
O arquivo \texttt{POSCAR} é um dos arquivos de entrada do VASP que contém informações sobre a estrutura
cristalina do sistema a ser simulado. Ele define a geometria da célula unitária, as posições dos
átomos e os tipos de átomos presentes na simulação. O formato do arquivo \texttt{POSCAR} é o seguinte:

\begin{minted}[linenos, frame=single, bgcolor=lightgray!20]{text}
EG-Graphene                             
   1.0000000000000000     
     3.8415129872762668   -0.0000076499719885    0.0000000000000000
    -1.9207631187581891    3.3268440109589794    0.0000000000000000
     0.0000000000000000    0.0000000000000000   15.0000000000000000
   C 
     4
Direct
     0.1210004546451131    0.2420030607848673    0.5114700199999973
     0.7579969542151304    0.8789995233548851    0.5114700199999973
     0.1209999839420490    0.8789999940579492    0.5114700199999973
     0.3333326271977057    0.6666674018022931    0.5114700199999973

  0.00000000E+00  0.00000000E+00  0.00000000E+00
  0.00000000E+00  0.00000000E+00  0.00000000E+00
  0.00000000E+00  0.00000000E+00  0.00000000E+00
  0.00000000E+00  0.00000000E+00  0.00000000E+00
\end{minted}

Na linha 1 do arquivo \texttt{POSCAR}, é especificado o nome do sistema ou uma descrição da
estrutura. A linha 2 contém o fator de escala, usado para dimensionar os vetores da rede. As
linhas 3 a 5 contêm os vetores da rede, que definem a geometria da célula unitária. A linha 6 lista
os tipos de átomos presentes na simulação, enquanto a linha 7 indica o número de átomos de cada
tipo. A linha 8 pode conter a palavra-chave \texttt{Selective dynamics} (opcional), que permite
especificar quais átomos podem se mover durante a simulação. A linha 9 indica se as posições dos
átomos estão em coordenadas diretas (frac) ou cartesianas (cart). As linhas subsequentes contêm as
posições dos átomos na célula unitária, de acordo com o número de átomos especificado na linha 7.

De posse dessas informações, vemos que o material estudado é o EG-Graphene, com célula unitária
definida pelos vetores da rede e contendo 4 átomos de carbono. As posições dos átomos são fornecidas
em coordenadas diretas, e a simulação é realizada em uma célula unitária com dimensões específicas.

\subsubsection{Arquivo \texttt{CONTCAR}}
O arquivo \texttt{CONTCAR} é gerado pelo VASP após a execução de uma simulação e contém as posições
finais dos átomos no sistema. O arquivo \texttt{CONTCAR} contém as informações sobre a estrutura,
por exemplo, as posições iônicas dos átomos. Seu formato é semelhante ao arquivo \texttt{POSCAR},
sendo escrito a cada passo iônico e após a conclusão do cálculo. Pode ser copiado para o arquivo
\texttt{POSCAR} para reiniciar a simulação a partir da última configuração.

No \textit{script} de execução do VASP no cluster, isso é feito de forma automática por meio de uma
instrução \texttt{while} que copia o arquivo \texttt{CONTCAR} para o arquivo \texttt{POSCAR} a cada
passo iônico, garantindo que a simulação continue a partir da última configuração obtida. Observe o
código abaixo:

\begin{minted}[linenos, frame=single, bgcolor=lightgray!20]{bash}
  i=50
  icount=0
  maxstep=10

  while [ $i -ne 1 ]
  do
      mpirun -hostfile $PBS_NODEFILE -np $NCPUS vasp_std > log6.out

      icount=$((icount+1))
      echo $icount >> log.dat

      if [ $icount -ge $maxstep ]; then
          exit 1
      fi

      cp CONTCAR POSCAR

      i=`grep F OSZICAR | tail -n1 | awk '{print $1;}'`
  done
\end{minted}

Observa-se que o código acima utiliza um loop \texttt{while} para executar o VASP até que a condição
de parada seja atendida. A cada iteração, o arquivo \texttt{CONTCAR} é copiado para o arquivo
\texttt{POSCAR}, garantindo que a simulação continue a partir da última configuração obtida.

